\documentclass[11pt]{article}
\usepackage[francais]{babel}

\setlength{\textwidth}{150mm}
\setlength{\textheight}{220mm}
\setlength{\hoffset}{-15mm}
\setlength{\voffset}{-20mm}


\title{Le Makefile de {\tt latmrg/xds}}
\author{Richard Simard}


\begin{document}
\maketitle
Pour que le Makefile fonctionne correctement, il est n\'ecessaire que 
les sous-r\'epertoires 
{\tt dirlilr, dirsilr, dirsisi, sparc20 ou linux }
existent, et que les programmes-ex\-em\-ples {\tt *.mod} 
soient dans leur r\'epertoire respectif {\tt dirxxxx}.

Les cibles possibles du Makefile sont:

\begin{itemize}
\item{\bf select} \\
c'est le premier programme qui doit exister avant de pouvoir 
reconstruire tous les modules et les programmes de {\tt latmrg}.
Il permet de choisir un des contextes {\tt LILR} possibles:
{\tt lilr, silr, sisi}. 

\item{\bf modules} \\
recompile tous les modules de base de {\tt latmrg}; ces fichiers
sont \'enum\'er\'es dans le fichier {\tt Latmrg.prj}. Toutefois,
les modules seront recompil\'es dans le contexte  courant de la variable
d'environnement {\bf LILR}. Par exemple, si on veut 
travailler uniquement dans le contexte {\bf sisi} (correspondant
au choix {\tt select si si}), alors on appellera d'abord {\bf sisi}
sur la ligne de commande du shell, ce qui nous placera dans ce contexte.
Alors toutes les recompilations des modules de base se feront 
dans ce contexte, et si on ne modifie qu'un seul fichier, seuls
les modules en d\'ependant seront recompil\'es.

\item{\bf modlilr} \\
force la recompilation de tous les modules de base dans le
contexte  {\tt lilr}.

\item{\bf modsilr} \\
force la recompilation de tous les modules de base dans le
contexte {\tt silr}.

\item{\bf modsisi} \\
force la recompilation de tous les modules de base dans le
contexte {\tt sisi}.

\item{\bf proglilr} \\
Recompile tous les programmes qui fonctionnent dans le contexte 
{\tt lilr}, mais en utilisant les fichiers {\tt *.o, *.sym, *.h}
de {\$XDSLIB}. %$
 On peut aussi recompiler un seul programme en 
le prenant comme cible, par exemple: {\tt make latl};
 mais dans ce cas, les fichiers  {\tt *.o, *.sym, *.h} de
{\tt latmrg/xds} seront utilis\'es, ce qui pr\'esume qu'on a
choisi le contexte appropri\'e en appelant {\tt lilr}. 
Ceci nous permet de tester {\bf nos propres fichiers  latmrg} 
\`a l'aide des diff\'erents programmes.

REMARQUE: si on veut compiler un programme utilisant les modules
de  {\tt latmrg} sans passer par le Makefile, alors il faut d'abord choisir
le contexte appropri\'e dans le shell. Par exemple,
 pour compiler un programme {\tt prog.mod}
 qui fonctionne dans le contexte de {\tt select si si}, il faudra
appeler {\tt sisi} sur la ligne de commande du shell, puis 
compiler avec \\
\hspace*{10mm}  {\tt xm =m prog.mod}.


\item{\bf progsilr} \\
Comme {\bf proglilr}, mais dans le contexte {\tt silr}.

\item{\bf progsisi} \\
Comme {\bf proglilr}, mais dans le contexte {\tt sisi}.


\item{\bf allprog} \\
Recompile tous les programmes de tous les contextes en utilisant 
les biblioth\`eques de \$XDSLIB. %%$

\item{\bf alllilr} \\
Recompile tous les modules de base, la biblioth\`eque, 
ainsi que tous les programmes 
qui fonctionnent dans le contexte {\tt lilr} en utilisant 
la biblioth\`eque et les fichiers de \$XDSLIB. %%$ 

\item{\bf allsilr} \\
Recompile tous les modules de base, la biblioth\`eque, 
 ainsi que tous les programmes 
qui fonctionnent dans le contexte {\tt silr} en utilisant 
la biblioth\`eque et les fichiers de \$XDSLIB. %%$

\item{\bf allsisi} \\
Recompile tous les modules de base, la biblioth\`eque, 
 ainsi que tous les programmes 
qui fonctionnent dans le contexte {\tt sisi} en utilisant 
la biblioth\`eque et les fichiers de \$XDSLIB. %%$

\item{\bf all}\\
 Identique \`a {\bf force}

\item{\bf force} \\
Force la reconstruction et la recompilation de tous les fichiers, les
modules, les biblioth\`eques et les programmes.

\item{\bf lib} \\
Reconstruit la biblioth\`eque liblat\$\{LILR\}99.a %$
selon le contexte du shell.

\item{\bf dellib} \\
D\'etruit toutes les biblioth\`eques de {\tt latmrg99}.

\item{\bf clean} \\
\'Elimine les fichiers de sauvegarde.

\item{\bf cleanprog} \\
\'Elimine les programmes ex\'ecutables ainsi que les 
fichiers de sauvegarde  de leur source.

\item{\bf cleanall} \\
\'Elimine les fichiers de sauvegarde, les ex\'ecutables et les 
fichiers auxiliaires de {\LaTeX}.

\end{itemize}

\end{document}On doit s'assurer
que les bons fichiers sont choisis.